\section{How we built our multi-agent research system}
\paperhead{How we built our multi-agent research system}{Jeremy Hadfield, Barry Zhang, Kenneth Lien, Florian Scholz, Jeremy Fox et al.}{Anthropic Engineering Blog, June 2025}{https://www.anthropic.com/engineering/multi-agent-research-system}{live page fetch, title present in page body}

\subsection{What the paper does}

Open-ended research cannot be planned in advance: the right next search depends on what the last one returned, so a fixed pipeline underfits the task. A single agent working in one context window is bounded twice over, by the context length it can hold and by the serial latency of its tool calls. This engineering report describes the production system behind Claude's Research feature and, unusually for an industry write-up, states the conditions under which its own architecture is the wrong choice.

The core idea is an orchestrator-worker split. A lead agent owns strategy, decomposition and synthesis; it spawns subagents that each search in their own private context window, compress what they find, and return a distilled result rather than raw evidence. Exploration is parallel and disposable, synthesis is central and singular. Attribution is factored out entirely into a separate \texttt{CitationAgent} that runs as a post-pass over the drafted report. The report is candid that this buys accuracy with tokens.

\subsection{How the method works}

The lead agent (Claude Opus 4) first thinks in extended thinking mode, used as a ``controllable scratchpad'', and persists its plan to external memory because context is truncated as it approaches 200{,}000 tokens. It then spawns subagents (Claude Sonnet 4), each handed four things: an objective, an output format, guidance on which tools and sources to use, and clear task boundaries. The report attributes duplicated work directly to underspecified task descriptions, giving an example where one subagent explored the 2021 automotive chip crisis while two others redundantly investigated 2025 supply chains. Subagents plan, call tools, then apply interleaved thinking to each tool result to judge quality, identify gaps and refine the next query, starting wide and narrowing down. Parallelism appears at two levels: the lead spawns subagents concurrently, and each subagent issues multiple tool calls concurrently. Effort is scaled by rules embedded in the prompt: one agent and 3 to 10 tool calls for simple fact-finding, 2 to 4 subagents at 10 to 15 calls each for direct comparisons, and 10 or more subagents with divided responsibilities for complex research. The lead synthesizes, then loops or exits. Finally \texttt{CitationAgent} processes the source documents together with the report to locate the specific spans that support each claim. Evaluation uses an LLM judge in a single call emitting a 0.0 to 1.0 score plus a pass/fail grade over factual accuracy, citation accuracy, completeness, source quality and tool efficiency, seeded with about 20 queries reflecting real usage; human testing caught failures the rubric missed, such as a bias toward SEO-optimized content over authoritative sources. Production practice includes full tracing of decision patterns without inspecting conversation contents, resumption of an agent from the point of failure, and rainbow deployments that shift traffic gradually so running agents are not disrupted.

\begin{center}
\begin{tikzpicture}[
  node distance=1.4cm, font=\scriptsize, >=Stealth, thick,
  bx/.style={draw, rounded corners=2pt, align=center, inner sep=3pt, minimum height=5.5mm}]
  \node[bx] (q) {user\\query};
  \node[bx, right=1.0cm of q] (lead) {LeadResearcher (Opus 4)\\plan / decompose / synthesize};
  \node[bx, right=1.0cm of lead] (cite) {CitationAgent\\post-pass binding};
  \node[bx, right=0.9cm of cite] (rep) {cited\\report};
  \node[bx, above=0.8cm of lead, xshift=-2.1cm] (s1) {Sub 1};
  \node[bx, above=0.8cm of lead] (s2) {Sub 2};
  \node[bx, above=0.8cm of lead, xshift=2.1cm] (s3) {Sub $n$};
  \node[bx, below=0.55cm of lead] (mem) {external memory (plan persisted near context limit)};
  \draw[->] (q) to (lead);
  \draw[->] (lead) to (cite);
  \draw[->] (cite) to (rep);
  \draw[<->] (lead) to (s1);
  \draw[<->] (lead) to (s2);
  \draw[<->] (lead) to (s3);
  \draw[<->] (lead) to (mem);
\end{tikzpicture}
\end{center}

\subsection{What it reports}

On an internal research evaluation, the multi-agent configuration with an Opus 4 lead and Sonnet 4 subagents outperformed a single-agent Opus 4 by 90.2\%, on the task of identifying board members across S\&P 500 Information Technology companies. In a BrowseComp analysis, three factors explained 95\% of variance, with token usage alone accounting for 80\%, followed by number of tool calls and model choice. Parallel tool calling cut research time by up to 90\% on complex queries, and improving tool descriptions produced a 40\% decrease in task completion time. The cost side is stated plainly: agents use roughly 4 times the tokens of a chat interaction and multi-agent systems roughly 15 times, so the pattern only pays where task value is high. The report also names its failure modes. The lead executes subagents synchronously and cannot steer them mid-execution, so one slow subagent blocks the system. Domains where all agents must share the same context or where subtasks carry many dependencies are called a poor fit, and the authors state that LLM agents are not yet good at coordinating and delegating to other agents in real time.

\subsection{How Mytheus uses this}
\begin{takeaway}
Cite \texttt{CitationAgent} as the industrial precedent for isolating attribution into its own model pass, then position the Mytheus \texttt{cite} stage as the stricter variant: theirs locates spans to attribute an already-written claim, whereas Mytheus drops any passage that cannot be bound to a resolvable \texttt{(meeting\_id, chunk\_id, speaker)} or \texttt{(document\_id, element\_id)} provenance tuple before the answer is assembled. Report the golden-set 100\% citation-resolution metric as the measurable difference. Their 15 times token multiple is the number the cost cascade should be argued against, but instrument it correctly: because token usage explained 80\% of performance variance in their analysis, the 86\% call-count reduction is not by itself a win, so measure tokens per confirmed decision or action against a single-chat baseline, not calls per meeting. Their spawn heuristics support the \texttt{PersonalListenerAgent} fan-out, which is genuinely breadth-first and independent per participant. They are a direct warning about \texttt{ContextGapAgent}: cross-participant queries are dependency-coupled by definition, exactly the case the report calls a poor fit, and a synchronous orchestrator with no mid-flight steering cannot hold a 1.5 s p95 peer round trip. Do not port the orchestrator-worker topology into \texttt{ContextGapAgent}. Keep the flat, deadline-carrying \texttt{agent-requests} / \texttt{agent-responses} exchange with \texttt{correlation\_id} and a hard timeout, and frame the fallback chain (peer agents, related-meeting agents, RAG, silence) as the bounded-time degradation path their blocking join lacks. Two operational patterns transfer directly: rainbow deployments map onto the 48-hour canary on 1\% of tenants, and their tracing of decision patterns without conversation contents is the template for observability that does not leak meeting text across the tenant boundary. Adopt their five judge criteria for the RAG golden set.
\end{takeaway}

\subsection{Limitations and what it leaves open}
\begin{itemize}[leftmargin=1.2em, itemsep=1pt, topsep=2pt]
  \item Every run produces one report for one requester; there is no notion of several simultaneous consumers with different roles, so nothing in the design addresses per-participant relevance scoring or role-conditioned output, which is the whole premise of \texttt{PersonalListenerAgent}.
  \item All work is batch retrieval over a static corpus reached through search tools: no streaming input, no events arriving mid-run, and the only speed figures given are aggregate wall-clock reductions rather than a bounded per-response deadline, so the architecture is never tested against a hard real-time budget.
  \item Coordination is strictly hierarchical and synchronous: subagents never query one another, and the lead cannot steer a subagent mid-execution, so the report supplies no evidence at all about peer-to-peer agent querying under a deadline.
  \item There is no multi-tenant story: no per-tenant token quota, no isolation of external memory or trace data, and no authorization model governing which agent may invoke which tool or which other agent.
  \item Outputs are returned as finished assertions with no calibrated confidence, no threshold at which a claim is withheld, and no human confirmation gate before a result is committed; the 90.2\% headline is also a single internal eval on one query family with no cost-matched single-agent baseline given the same token budget.
\end{itemize}
